% Sections I-V. Every number here is taken from a committed file; the source is named in a comment.

\section{Introduction}
Estimating a physical parameter from a few sensor readings usually means fitting a simulator to the data:
an optimiser proposes a parameter value, the simulator is run, and the mismatch with the readings decides the
next proposal. When each run is expensive, the number of runs is the cost. A common idea is to train a neural
network on simulated examples so that it predicts the parameter directly, and to use that prediction as the
starting point of the fit. The fit then keeps its accuracy and, one hopes, needs fewer runs.

We test this idea on a deliberately small problem where every part can be checked: recovering the viscosity
$\nu$ of the one-dimensional viscous Burgers equation from 4, 8 or 16 sensors read at four times, with a
known initial state and 0\%, 1\% or 5\% noise. A two-layer perceptron predicts $\log\nu$; a bounded
least-squares fit refines it. Our project proposal set a target: the learned start should halve the median
online cost at no loss of accuracy. We state that as a hypothesis and test it.

Three choices shape the study. First, the comparison includes a direct bounded scalar search (Brent's method),
which needs no starting point at all, and a nearest-neighbour start; a learned start that only beats a poor
cold start has shown little. Second, cost counts every simulator call, including those used for
finite-difference derivatives, and the offline cost of the network (data generation, tuning, training) is
divided by the saving per query to give a break-even number of queries. Third, the test-set evaluation was
written down and committed before the test set was used, and it is reported whatever the outcome.

The outcome is mostly negative. The learned start makes the least-squares fit cheaper than a midpoint start
in every setting we ran. Against Brent's search, it is cheaper at one stopping tolerance and not at a tighter
one, and the 2$\times$ target is not met. The remaining sections show why: the estimates already reach the
Cram\'er--Rao bound, so accuracy cannot separate the methods, and a one-dimensional search needs so few
simulator runs that little is left to save.

\section{Related Work}
\textbf{Learned warm starts.} Sambharya et al.~\cite{sambharya2024} train a network that maps the parameters
of a convex problem to a warm start for a fixed-point solver, and compare it with a cold start and a
nearest-neighbour warm start; they report iterations and solve times, and state training time separately
without setting it against the savings. Zhou et al.~\cite{zhou2023} predict the shape of a scattering
obstacle with a convolutional network and refine it by Gauss--Newton; their aim is to avoid local minima at
high frequency, they report running time per obstacle, and training time is excluded from that comparison.
Borisyak et al.~\cite{borisyak2023} predict the parameters of ODE growth models and optionally refine them by
gradient-based fitting. Oh et al.~\cite{oh2026} warm-start Newton solvers for nonlinear PDEs with a neural
operator and show that a low prediction error does not by itself give a usable start. Amos~\cite{amos2023}
surveys this family as amortised optimisation; in his terms our method is semi-amortised and trained by
regression.

\textbf{Amortised and surrogate-based parameter estimation.} Simulation-based inference replaces the
optimiser by a network trained on simulations~\cite{cranmer2020}; Tatsuoka et al.~\cite{tatsuoka2025} do so
with two fidelity levels to limit expensive forward runs. Roy Sarkar et al.~\cite{roysarkar2024} identify PDE
parameters from sparse, noisy data, Burgers among the examples, with a pretrained operator network in place
of the simulator, and Long et al.~\cite{long2025} train a recurrent network together with a numerical scheme
for each problem. We keep the simulator in the loop and learn only the starting point.

\textbf{Sensors and identifiability.} Neitzel et al.~\cite{neitzel2019} choose sensor locations by the
covariance of the least-squares estimator. We use the same Fisher-information view, not to place sensors but
to check whether our estimates use all the information the given sensors carry.

\textbf{Position of this work.} The mechanism is established; we do not claim it. What the two closest
works~\cite{sambharya2024,zhou2023} leave open, and what we measure on one small problem, is the comparison
against a direct scalar search, the break-even count once offline cost is included, and the dependence of the
saving on the stopping tolerance.

\section{Problem Setup}
\subsection{Equation and simulator}
We use the Burgers equation in the PDEBench convention~\cite{takamoto2022}, % configs/simulator_v1.json
\begin{equation}
u_t + \left(\tfrac{1}{2}u^2\right)_x = \frac{\nu}{\pi}\,u_{xx}, \qquad x\in[0,1)\ \text{periodic},
\end{equation}
with $\nu\in[0.01,0.1]$ and the initial state
$u_0(x)=\sum_{k=1}^{3} a_k \sin(2\pi k x+\phi_k)$, $a_k\sim U(0,1)$, $\phi_k\sim U(0,2\pi)$. The amplitudes and
phases are known to every method; only $\nu$ is unknown. The solver is our own NumPy implementation of the
finite-volume scheme used by PDEBench (MUSCL reconstruction with its limiter, Rusanov flux, central
diffusive flux, two-stage midpoint time stepping, CFL 0.25) on 256 cells in double precision. No PDEBench
code is included.

The solver passes eleven tests, % tests/test_simulator.py
among them convergence to the exact Cole--Hopf solution, the small-amplitude heat limit with diffusivity
$\nu/\pi$, mass conservation and periodic shift invariance. Its cost depends on $\nu$ because the explicit
diffusion step limits the time step: one run to $t=1$ takes 1{,}686 steps at $\nu=0.01$ and 16{,}692 at
$\nu=0.1$. % RESEARCH_LOG.md, "Cost observation"
A method's cost therefore depends on which values of $\nu$ it tries, not only on how many.

\subsection{Observations}
Sixteen sensor slots sit at $x_s=(s+0.5)/16$ and are read at $t\in\{0.25,0.5,0.75,1\}$. A case uses 16, 8 or
4 of them (every slot, every second, every fourth). Gaussian noise of standard deviation
$\sigma=\ell\cdot\mathrm{RMS}$ is added, where $\ell\in\{0,0.01,0.05\}$ and RMS is taken over the clean
$16\times4$ readings of that trajectory. % blinit/data.py make_observations
Each trajectory thus gives nine cases (three sensor counts $\times$ three noise levels).

\subsection{Data}
No public dataset is used; all data are generated by the repository's solver from versioned configuration
files and regenerate identically. % configs/phaseB_v2.json, data/phaseB_v2/integrity.json
There are 3{,}000 training, 50 validation and 50 test trajectories (27{,}000 / 450 / 450 cases), with $\nu$
drawn log-uniformly. Splits are made by the seed of the initial state, so all noise and sensor variants of a
trajectory stay in one split. Nine automated checks run at build time, covering seed overlap between splits,
normalisation statistics taken from training data only, and the mapping from readings to sensor positions.

\section{Methods}
All methods estimate $\theta=\log\nu$ within $[\log 0.01,\log 0.1]$ by minimising
$J(\theta)=\tfrac12\lVert f(\theta)-y\rVert^2$ over the available readings, where $f$ is the simulator
followed by interpolation to the sensors. They share one problem object that counts every simulator call.
% blinit/inverse.py

\textbf{Brent.} Bounded scalar minimisation of $J$ (SciPy \texttt{minimize\_scalar}, method
\texttt{bounded}) with absolute tolerance \texttt{xatol} on $\theta$. It needs no starting point.

\textbf{Start + LSQ.} A starting value $\theta_0$ followed by bounded least squares (SciPy
\texttt{least\_squares}, trust-region reflective, two-point finite differences with relative step
$10^{-6}$, tolerance \texttt{xtol}). Four starts are compared:
\emph{midpoint} of the interval; \emph{grid}, the best of five log-spaced values (five extra simulator runs,
counted); \emph{kNN}, the inverse-distance-weighted mean of the $k$ nearest training cases with the same
sensor pattern; and \emph{MLP}.

\textbf{MLP.} A perceptron with two hidden layers of 64 ReLU units (scikit-learn \texttt{MLPRegressor},
Adam, batch size 64) maps 134 inputs to $\log\nu$: the 64 standardised readings with missing ones set to
zero, 64 availability flags, and the six numbers $a_k\cos\phi_k$, $a_k\sin\phi_k$. % configs/phaseC_v2.json
Training stops early on the validation loss. Four settings of learning rate and weight decay were tried and
the one with the lowest validation loss kept; three seeds were then trained. % models/phaseC_v2/tuning.json
kNN's $k$ and feature weight were chosen on validation as well. The test split was not loaded during any of
this.

\textbf{Cost.} The cost of a query is its number of solver steps times 144.3\,\textmu s, the median time per
step measured in sequential runs. % RESEARCH_LOG.md, Phase A run 2
We use steps rather than wall-clock time because experiments ran on parallel workers of a laptop whose
timings were not clean. The offline cost of the MLP is the wall time of training-data generation multiplied
by the number of workers (an upper bound), plus tuning and training. % blinit/experiments/analyze_d.py
The break-even count is this offline cost divided by the mean saving per query against the cheapest
classical method, with the MLP's inference time subtracted from the saving.

\section{Evaluation Protocol}
\textbf{Tolerances.} Stopping tolerances were fixed on development data: Brent \texttt{xatol} $10^{-3}$ and
LSQ \texttt{xtol} $10^{-3}$ (primary). Because the LSQ rule can stop looser than Brent in the worst case, a
second setting with LSQ \texttt{xtol} $10^{-4}$ (accuracy-matched) was added before testing.

\textbf{Pre-registration.} Before the test split was used, a document committed to the
repository % docs/preregistration_pilot_test.md
fixed the models, the methods, both tolerance settings, the cost metric, the statistic (median paired
difference in CPU-seconds between a learned method and Brent on the same case, with a 95\% bootstrap interval
that resamples trajectories), and the rule for success: the saving counts as positive only if the interval
excludes zero under both tolerance settings for at least two of the three MLP seeds. The 2$\times$ target is
judged on the ratio of median costs. The test split was evaluated once per setting and nothing was changed
afterwards. The robustness sets of Section~\ref{sec:robust} were pre-registered in the same way before they
were generated. % docs/preregistration_phaseF.md

\textbf{What validation was used for.} Early stopping, the choice among four MLP settings, the kNN settings,
the tolerance curves of Section~\ref{sec:tol} and the ablations. Validation numbers are labelled as such.
